\notechapter{N-20221211}{Derivatives: Definition, Rules, Elementary Functions, and Inverse Trigonometric Derivatives}

\section{Average velocity and instantaneous velocity}

The first page starts from a graph and the idea of velocity.  The average rate of change of $f$ on $[x_0,x_0+\Delta x]$ is
\[
  \frac{f(x_0+\Delta x)-f(x_0)}{\Delta x}.
\]
As $\Delta x\to0$, the secant line approaches the tangent line if the limit exists.  The instantaneous rate of change is the derivative:
\[
  f'(x_0)=\lim_{\Delta x\to0}
  \frac{f(x_0+\Delta x)-f(x_0)}{\Delta x}.
\]
Geometrically, this is the slope of the tangent line.

\section{Derivative notation}

The note lists the notations
\[
  f'(x),\qquad y',\qquad \frac{dy}{dx},\qquad \frac{df}{dx}.
\]
It also distinguishes the derivative at a point from the derivative function.  The derivative at $x_0$ is a number; the derivative function assigns such a number to each point where the derivative exists.

The tangent line formula is
\[
  y-f(x_0)=f'(x_0)(x-x_0).
\]

\section{Power functions}

For positive integer $n$,
\[
  (x^n)'=nx^{n-1}.
\]
This follows from
\[
  (x+h)^n-x^n
  =\sum_{k=1}^n \binom nk x^{n-k}h^k,
\]
and after dividing by $h$ and taking $h\to0$, only the $k=1$ term remains.

The note extends the formula to rational powers by implicit differentiation.  If
\[
  y=x^{m/n},
\]
then
\[
  y^n=x^m.
\]
Differentiating gives
\[
  ny^{n-1}y'=mx^{m-1},
\]
so
\[
  y'=\frac mn x^{m/n-1}.
\]

\section{Exponential and logarithmic derivatives}

The page records
\[
  (a^x)'=a^x\ln a,
  \qquad
  (e^x)'=e^x.
\]
For logarithms,
\[
  (\ln x)'=\frac1x,
  \qquad
  (\log_a x)'=\frac1{x\ln a}.
\]
A useful derivation uses inverse functions: since $y=\ln x$ means $x=e^y$, differentiating gives
\[
  1=e^y y'=xy',
\]
so $y'=1/x$.

\section{Trigonometric derivatives}

Using the standard limits and addition formulas,
\[
  (\sin x)'=\cos x,\qquad
  (\cos x)'=-\sin x.
\]
Then quotient rule gives
\[
  (\tan x)'=\sec^2x,\qquad
  (\cot x)'=-\csc^2x.
\]
The reciprocal functions satisfy
\[
  (\sec x)'=\sec x\tan x,\qquad
  (\csc x)'=-\csc x\cot x.
\]

\section{Inverse trigonometric derivatives}

For $y=\arcsin x$,
\[
  x=\sin y
\]
and so
\[
  1=\cos y\,y'.
\]
Since $\cos y=\sqrt{1-x^2}$ on the principal branch,
\[
  (\arcsin x)'=\frac1{\sqrt{1-x^2}}.
\]
Similarly,
\[
  (\arccos x)'=-\frac1{\sqrt{1-x^2}},\qquad
  (\arctan x)'=\frac1{1+x^2}.
\]
The note also derives inverse secant and inverse cosecant derivatives by drawing a right triangle and expressing the missing side length in terms of $x$.

\section{Examples}

A typical example differentiates
\[
  y=\sqrt{x^2+1}.
\]
By chain rule,
\[
  y'=\frac{x}{\sqrt{x^2+1}}.
\]
Another example differentiates a quotient of trigonometric functions.  The handwritten solution usually first rewrites the expression using identities, then differentiates.  This is often simpler than applying quotient rule blindly.

\section{Context and outlook}

This note builds the derivative table from a few principles: definition by difference quotient, binomial expansion for powers, inverse-function differentiation for logarithms and inverse trigonometric functions, and quotient/chain rules for combinations.  The goal is not memorizing isolated formulas but recognizing which rule creates each formula.

\section{Derivative as first-order linearization}

The derivative at \(a\) is best read as a first-order approximation:
\[
  f(a+h)=f(a)+f'(a)h+o(h).
\]
The number \(f'(a)\) is the linear map that best approximates the change of \(f\) near \(a\).  In one dimension, linear maps are multiplication by a number, so this viewpoint collapses to the familiar tangent slope.

This explains why the chain rule has the form it does.  If
\[
  f(a+h)=f(a)+f'(a)h+o(h),\qquad
  g(b+k)=g(b)+g'(b)k+o(k),\quad b=f(a),
\]
then
\[
  g(f(a+h))=g(f(a))+g'(f(a))f'(a)h+o(h).
\]
Thus
\[
  (g\circ f)'(a)=g'(f(a))f'(a).
\]
The chain rule is composition of linear approximations.

\section{Inverse trigonometric derivatives and the inverse function theorem}

The inverse-function formula
\[
  (f^{-1})'(y)=\frac{1}{f'(x)},\qquad y=f(x),
\]
is a one-dimensional form of the inverse function theorem.  If \(f\) is continuously differentiable near \(x_0\) and \(f'(x_0)\ne0\), then \(f\) has a local differentiable inverse near \(f(x_0)\).

For \(y=\arcsin x\), the equation \(x=\sin y\) is inverted only after restricting sine to
\[
  [-\pi/2,\pi/2].
\]
On this branch, \(\cos y\ge0\), so
\[
  1=\cos y\frac{dy}{dx},\qquad
  \frac{dy}{dx}=\frac{1}{\cos y}=\frac{1}{\sqrt{1-x^2}}.
\]
The square root sign is not arbitrary; it comes from the chosen branch of the inverse.

\section{Endpoint blow-up and failure of local invertibility}

The derivative
\[
  (\arcsin x)'=\frac{1}{\sqrt{1-x^2}}
\]
blows up as \(x\to\pm1\).  Geometrically, near \(y=\pm\pi/2\), the graph of \(\sin y\) has horizontal tangent:
\[
  \cos y=0.
\]
The inverse function theorem fails exactly where \(f'(y)=0\).  The inverse may still exist globally on a restricted interval, but it is no longer differentiable with a finite derivative at the endpoint.

This is the conceptual reason inverse-trigonometric formulas have domain restrictions and singular denominators.

\section{Logarithm as inverse and as group homomorphism}

The derivative of \(\ln x\) follows from \(x=e^y\):
\[
  1=e^y y'=xy',\qquad y'=\frac1x.
\]
But \(\ln\) also has a structural meaning:
\[
  \ln(xy)=\ln x+\ln y.
\]
It converts multiplication in \((0,\infty)\) into addition in \(\mathbb R\).  Its derivative \(1/x\) is the infinitesimal version of multiplicative scaling.

In Lie-group language, \((0,\infty)\) is a one-dimensional Lie group under multiplication, and \(\ln\) identifies it with the additive group \(\mathbb R\).  The elementary derivative formula is the local shadow of this group isomorphism.

\section{Automatic differentiation and computation graphs}

The product, quotient, and chain rules are the mathematical basis of automatic differentiation.  A complicated expression is a computation graph.  Each node knows its local derivative, and the chain rule propagates derivative information through the graph.

For example, if
\[
  y=\sqrt{x^2+1},\qquad u=x^2+1,\qquad y=\sqrt u,
\]
then
\[
  \frac{dy}{dx}=\frac{dy}{du}\frac{du}{dx}=\frac{1}{2\sqrt u}\cdot2x=\frac{x}{\sqrt{x^2+1}}.
\]
The ordinary chain-rule computation is already reverse or forward mode automatic differentiation in miniature.

\section{Jets and higher-order local data}

The \(k\)-jet of a function at \(a\) records its derivatives up to order \(k\):
\[
  j_a^k f=(f(a),f'(a),\ldots,f^{(k)}(a)).
\]
The Taylor polynomial is the coordinate expression of this finite local data:
\[
  T_kf(a+h)=\sum_{j=0}^k \frac{f^{(j)}(a)}{j!}h^j.
\]
The first derivative is therefore not an isolated operation; it is the first nontrivial layer of a tower of local approximations.

\section{Differentials and one-forms}

The notation \(dy/dx\) is sometimes taught as a fraction-like symbol.  A safer advanced interpretation is through differentials.  For \(y=f(x)\),
\[
  dy=f'(x)\,dx.
\]
Here \(dx\) is the basic one-form on the line, and \(dy\) is its pullback through \(f\).  The chain rule becomes the rule for pulling back one-forms:
\[
  (g\circ f)^*(dz)=f^*(g^*(dz)).
\]
This gives a structural meaning to the notation without treating infinitesimals as informal numbers.

\section{Derivatives as linearization and pullback}

The derivative table is built from a few structural principles: linear approximation, inverse functions, branch choices, composition, and local Taylor data.  The formulas for powers, logarithms, trigonometric functions, and inverse trigonometric functions are not separate facts; they are examples of how local linear behavior is transported through algebraic operations and inverses.
